\section{The library: archetypes, entries, and their evidence}
\label{app:typologies}

This appendix reproduces the full library at methodology version
\methver{}, as it appears in \texttt{config/nbs\_typologies.yaml}. All
temperature values are daytime, pedestrian-level \emph{air} temperature
reductions in degrees Celsius, relative to an equivalent unimproved site.

The library has two levels (\Cref{sec:archetypes}). \Cref{tab:app-archetypes,%
tab:app-cobenefits,tab:app-citations,tab:app-notes} give the \num{18} cooling
archetypes, which carry every performance value and every citation.
\Cref{tab:app-entries} lists the \num{121} catalogue entries and the archetype
each one inherits; an entry carries no performance value of its own.
\Cref{tab:app-overrides} lists the \num{21} per-entry suitability overrides.

\subsection{Archetype values and inherited suitability conditions}

{\small
\setlength{\tabcolsep}{4pt}
\begin{longtable}{@{}lrcccrl@{}}
\caption[Archetype values and suitability conditions]{Adopted values and
suitability conditions for all \num{18} cooling archetypes. \emph{Bldg.\
energy} indicates whether the archetype is flagged as capable of affecting
adjacent building cooling loads, a precondition for the energy derivation of
\Cref{sec:energy}. The minimum area and the soil and irrigation requirements
are inherited by every catalogue entry resolving to the archetype, subject to
the overrides of \Cref{tab:app-overrides}.}
\label{tab:app-archetypes}\\
\toprule
\textbf{Archetype} & \textbf{Base} & \textbf{Envelope} &
\textbf{Conf.} & \textbf{Bldg.} & \textbf{Min.\ area} & \textbf{Soil / water} \\
& \textbf{score} & (\si{\degreeCelsius}) & & \textbf{energy} &
(\si{\metre\squared}) & \textbf{required} \\
\midrule
\endfirsthead
\multicolumn{7}{@{}l}{\small\itshape \Cref{tab:app-archetypes} continued from
previous page}\\
\toprule
\textbf{Archetype} & \textbf{Base} & \textbf{Envelope} &
\textbf{Conf.} & \textbf{Bldg.} & \textbf{Min.\ area} & \textbf{Soil / water} \\
\midrule
\endhead
\midrule
\multicolumn{7}{r@{}}{\small\itshape continued on next page}\\
\endfoot
\bottomrule
\endlastfoot
\multicolumn{7}{@{}l}{\textit{Green}}\\
Street tree canopy & 75 & 0.5--3.0 & High & Yes & 100 & Limited / occas. \\
Continuous shaded corridor & 80 & 0.5--3.0 & Medium & No & 200 & Limited / occas. \\
Small green area with trees & 65 & 0.3--1.5 & High & No & 200 & Limited / occas. \\
Dense tree canopy & 90 & 1.0--3.0 & High & Yes & 2000 & Moderate / occas. \\
Established park & 70 & 0.5--2.0 & High & No & 500 & Limited / occas. \\
Non-canopy vegetation & 50 & 0.0--1.5 & Medium & No & 50 & Limited / occas. \\
Productive tree canopy & 60 & 0.5--2.0 & Low & No & 500 & Moderate / reliable \\
Productive ground-level planting & 45 & 0.0--1.5 & Low & No & 200 & Moderate / reliable \\
\addlinespace
\multicolumn{7}{@{}l}{\textit{Building}}\\
Extensive green roof & 45 & 0.1--1.0 & Medium & Yes & 50 & None / occas. \\
Green facade or living wall & 50 & 0.3--2.0 & Low & Yes & 20 & None / reliable \\
\addlinespace
\multicolumn{7}{@{}l}{\textit{Blue-green}}\\
Bioretention and drainage planting & 55 & 0.1--0.8 & Low & No & 50 & Moderate / none \\
Blue-green corridor & 85 & 1.0--3.0 & Medium & No & 1000 & Moderate / none \\
Riparian restoration & 85 & 1.0--3.0 & Medium & No & 1000 & Moderate / none \\
Large open water body & 75 & 0.1--3.0 & Medium & No & 1000 & Moderate / none \\
Small constructed water feature & 40 & 0.0--1.0 & Medium & No & 50 & Moderate / none \\
\addlinespace
\multicolumn{7}{@{}l}{\textit{Hybrid}}\\
Permeable shaded hardscape & 70 & 0.5--2.5 & Medium & No & 200 & Limited / occas. \\
Enclosed courtyard greening & 60 & 0.3--1.5 & Low & Yes & 50 & Limited / occas. \\
Vegetated shade structure & 55 & 0.0--2.5 & Low & No & 20 & None / reliable \\
\addlinespace
\end{longtable}
}

No archetype declares an unsuitable climate zone at version \methver{}. Climate
influences performance through the adjustment matrix of
\Cref{tab:climate-matrix} rather than through outright disqualification.

\subsection{Co-benefit defaults}

The primary cooling mechanism is declared \emph{per catalogue entry} rather
than per archetype, because it is the one descriptive attribute the source
catalogue supplies at entry resolution; it is reported in the result and enters
no formula. The co-benefit defaults below are archetype properties, inherited
by every entry, and are unioned across components in a package
(\Cref{tab:package-rules}).

\begin{longtable}{@{}>{\raggedright\arraybackslash}p{62mm}cccc@{}}
\caption[Co-benefit defaults]{Default co-benefit levels per archetype. The
library supplies no default for the \emph{social inclusion} sub-indicator of
\Cref{eq:cobenefit}; absent a user override it falls to the neutral
\emph{unknown} value and is itemised as an applied assumption.}
\label{tab:app-cobenefits}\\
\toprule
\textbf{Archetype} & \textbf{Biodiv.} &
\textbf{Storm.} & \textbf{Health} & \textbf{Urban q.} \\
\midrule
\endfirsthead
\toprule
\textbf{Archetype} & \textbf{Biodiv.} &
\textbf{Storm.} & \textbf{Health} & \textbf{Urban q.} \\
\midrule
\endhead
\bottomrule
\endlastfoot
Street tree canopy & Med & Med & High & High \\
Continuous shaded corridor & Med & Low & High & High \\
Small green area with trees & Med & Med & High & High \\
Dense tree canopy & High & High & High & High \\
Established park & High & Med & High & High \\
Extensive green roof & Med & High & Low & Med \\
Green facade or living wall & Low & Low & Low & Med \\
Bioretention and drainage planting & Med & High & Low & Med \\
Blue-green corridor & High & High & High & High \\
Riparian restoration & High & High & Med & High \\
Permeable shaded hardscape & Low & High & Med & High \\
Enclosed courtyard greening & Med & Med & Med & High \\
Large open water body & High & High & Med & High \\
Small constructed water feature & Med & High & Low & Med \\
Non-canopy vegetation & Med & Med & Low & Med \\
Vegetated shade structure & Low & Low & Med & Med \\
Productive tree canopy & Med & Med & High & High \\
Productive ground-level planting & Med & Low & High & Med \\
\end{longtable}

\subsection{Citations and the findings they support}

Every performance value in the library carries at least one citation, and each
citation records the specific finding it supports. Configuration loading fails
if a resolved entry is uncited or if a citation names a source absent from the
bibliography (\Cref{sec:evidence-rules}). The findings below are reproduced
from the configuration.

\begin{longtable}{@{}>{\raggedright\arraybackslash}p{30mm} l
                    >{\raggedright\arraybackslash}p{72mm}@{}}
\caption[Archetype citations and findings]{Citations supporting each
archetype's adopted values, with the finding each citation supports. Metrics
are stated where a source reports something other than air temperature.}
\label{tab:app-citations}\\
\toprule
\textbf{Archetype} & \textbf{Source} & \textbf{Finding relied upon} \\
\midrule
\endfirsthead
\multicolumn{3}{@{}l}{\small\itshape \Cref{tab:app-citations} continued}\\
\toprule
\textbf{Archetype} & \textbf{Source} & \textbf{Finding relied upon} \\
\midrule
\endhead
\midrule
\multicolumn{3}{r@{}}{\small\itshape continued on next page}\\
\endfoot
\bottomrule
\endlastfoot
Street tree canopy & \texttt{keravec2026} &
Street trees: 1.5\,\si{\degreeCelsius} central estimate for pedestrian-level daytime air temperature; 4.0 $\pm$ 1.6\,\si{\degreeCelsius} in tropical (A) and 1.7 $\pm$ 1.4\,\si{\degreeCelsius} in arid (B) climates. \\
 & \texttt{ziter2019} &
Air temperature decreases nonlinearly with canopy cover, greatest cooling above $\sim$40\% cover; peak effect at 60--90 m scale. \\
 & \texttt{bowler2010} &
Empirical studies broadly support cooler air beneath tree canopy. \\
 & \texttt{rahman2020} &
Species and leaf area index materially change delivered cooling. \\
 & \texttt{kumar2024} &
Street trees measured at up to 2.8\,\si{\degreeCelsius} by in-situ monitoring, inside this envelope: newer and broader evidence supports the conservative calibration rather than overturning it. \\
\addlinespace
Continuous shaded corridor & \texttt{keravec2026} &
Street trees 1.5\,\si{\degreeCelsius} and shading devices 2.0\,\si{\degreeCelsius} pedestrian-level daytime air temperature reduction. \\
 & \texttt{zolch2016} &
Shade placement in heat-exposed locations outperforms maximising total green cover; trees -13\% PET. \\
\addlinespace
Small green area with trees & \texttt{bowler2010} &
Meta-analysis: park on average 0.94\,\si{\degreeCelsius} cooler by day; the review reports that larger parks and parks containing trees tended to be cooler, stated qualitatively. \\
 & \texttt{keravec2026} &
Parks 1.3\,\si{\degreeCelsius} pedestrian-level daytime air temperature reduction. \\
\addlinespace
Dense tree canopy & \texttt{keravec2026} &
Urban forests 2.1\,\si{\degreeCelsius} pedestrian-level daytime air temperature reduction, among the highest-performing green measures. \\
 & \texttt{ziter2019} &
Greatest cooling where canopy cover exceeds 40\%; effect grows with patch scale to 60--90 m. \\
 & \texttt{bowler2010} &
Larger, tree-rich green sites cooler during the day. \\
\addlinespace
Established park & \texttt{bowler2010} &
Park mean 0.94\,\si{\degreeCelsius} daytime cooling; parks with trees reported cooler than open turf. \\
 & \texttt{keravec2026} &
Parks 1.3\,\si{\degreeCelsius}, green areas 2.0\,\si{\degreeCelsius} pedestrian-level daytime air temperature reduction. \\
 & \texttt{kraemer2022} &
Under 2018--2019 drought and heat, maximum spatially averaged park-to-built cooling was 1.1\,\si{\degreeCelsius} in the morning, falling below 1\,\si{\degreeCelsius} in the afternoon and confined to shaded areas. \\
\addlinespace
Extensive green roof & \texttt{keravec2026} &
Green roofs 1.1\,\si{\degreeCelsius} pedestrian-level daytime air temperature reduction, lowest-performing green measure in the synthesis. \\
 & \texttt{zolch2016} &
Green roof effects at pedestrian level negligible. \\
 & \texttt{santamouris2014} &
0.3--3\,\si{\kelvin} average urban ambient reduction, but only for city-scale deployment in simulation studies, not single-building installation. \\
\addlinespace
Green facade or living wall & \texttt{keravec2026} &
Green walls 3.0\,\si{\degreeCelsius} central, with fivefold spread across climate subzones (Cwa 5.7, BWh 5.5, Csa 2.9, Cfa/Cfb 1.3, Af 1.0\,\si{\degreeCelsius}). \\
 & \texttt{zolch2016} &
Green facades -5\% to -10\% PET, below trees at -13\%. \\
\addlinespace
Bioretention and drainage planting & \texttt{keravec2026} &
No bioswale-specific estimate available; neighbouring categories (green areas 2.0\,\si{\degreeCelsius}, water bodies 2.1\,\si{\degreeCelsius}) are not transferable to small drainage features. \\
 & \texttt{gunawardena2017} &
Tree-dominated greenspace delivers heat-stress relief; small blue features are not identified as significant coolers. \\
\addlinespace
Blue-green corridor & \texttt{keravec2026} &
Water bodies 2.1\,\si{\degreeCelsius} and urban forests 2.1\,\si{\degreeCelsius} pedestrian-level daytime air temperature reduction; water bodies and parks reported as largely climate-independent. \\
 & \texttt{gunawardena2017} &
Green and blue features together offer synergistic benefits; poorly designed bluespace may exacerbate heat stress in humid conditions. \\
\addlinespace
Riparian restoration & \texttt{keravec2026} &
Water bodies 2.1\,\si{\degreeCelsius}, climate-robust; urban forests 2.1\,\si{\degreeCelsius} pedestrian-level daytime air temperature reduction. \\
 & \texttt{gunawardena2017} &
Bluespace cooling mechanisms; humid-condition heat-stress caveat. \\
\addlinespace
Permeable shaded hardscape & \texttt{keravec2026} &
Permeable pavements 2.9\,\si{\degreeCelsius}, cool pavements 2.0\,\si{\degreeCelsius}, shading devices 2.0\,\si{\degreeCelsius} pedestrian-level daytime air temperature reduction. \\
 & \texttt{zolch2016} &
Shade positioned in heat-exposed locations is the decisive design factor. \\
\addlinespace
Enclosed courtyard greening & \texttt{keravec2026} &
Green areas 2.0\,\si{\degreeCelsius} and green walls 3.0\,\si{\degreeCelsius}; enclosed-courtyard geometry not separately resolved. \\
 & \texttt{ziter2019} &
Cooling scales with patch size; courtyards fall below the 60--90 m optimum. \\
 & \texttt{zolch2016} &
Micro-scale greening in dense residential fabric reduces PET; placement dominates quantity. \\
\addlinespace
Large open water body & \texttt{keravec2026} &
Water bodies 2.1\,\si{\degreeCelsius} pedestrian-level daytime air temperature reduction, reported as largely climate-independent. \\
 & \texttt{volker2013} &
Meta-analysis of 27 studies: pooled cooling of 2.5\,\si{\kelvin} (95\% CI 1.9--3.2\,\si{\kelvin}) in the warmest months, attributed to urban blue sites WHEN REMOTE-SENSING DATA ARE INCLUDED. A single pooled estimate across ponds, lakes and rivers together; it does not differentiate cooling by water-body size. \\
 & \texttt{yao2023a} &
Year-long field measurement of urban lakes in Nanjing and Guangzhou: weak daytime urban cool island intensity of 0.1--0.6\,\si{\degreeCelsius}, against a 1.2--1.3\,\si{\degreeCelsius} heat island intensity at night in the warm months. \\
 & \texttt{ampatzidis2020} &
A clear disparity exists between cooling potentials reported by remote sensing and those from field measurement or numerical simulation; at night blue spaces may exacerbate the urban heat island. \\
\addlinespace
Small constructed water feature & \texttt{jacobs2020} &
Sixteen simulated urban water bodies: afternoon air temperature in surrounding spaces reduced by typically 0.2\,\si{\degreeCelsius}, maximum 0.6\,\si{\degreeCelsius}. Of 20 papers synthesised for water bodies without fountains, 14 (70\%) report an air-temperature effect of 1\,\si{\degreeCelsius} or less; local thermal effects of small water bodies can be considered negligible in design practice. \\
 & \texttt{yao2023b} &
Field measurement of an urban pond: moderate daytime cooling of 0.6\,\si{\degreeCelsius} against a reference soil site, and a pronounced nocturnal warming effect of 1.8\,\si{\degreeCelsius}. \\
 & \texttt{ampatzidis2020} &
Blue-space cooling varies with time of day, season and location; at night blue spaces may exacerbate the urban heat island, and sensible cooling can be offset by increased water-vapour content. \\
\addlinespace
Non-canopy vegetation & \texttt{armson2012} &
Measured on the same plots: grass reduced maximum SURFACE temperature by up to 24\,\si{\degreeCelsius} but had little effect upon GLOBE temperature, while shading reduced globe temperature by 5--7\,\si{\degreeCelsius}. Grass has little effect on local air or globe temperature; tree shade provides effective local cooling. \\
 & \texttt{gill2013} &
Modelled: grass cooling can be reduced or lost in summer droughts when soils dry out, an effect projected to become more pronounced under climate change. \\
 & \texttt{kraemer2022} &
Field measurement under drought: the grass-dominated park was the hottest area during the daytime, while the tree-dominated park was consistently cooler. \\
 & \texttt{bowler2010} &
Parks containing trees reported cooler than open turf, stated qualitatively. \\
 & \texttt{kumar2024} &
Typology-resolved review placing grass and herbaceous systems well below canopy systems for air-temperature cooling. \\
\addlinespace
Vegetated shade structure & \texttt{chafer2020} &
The only measured comparison of a vine-covered pergola against an identical bare support: maximum daytime air temperature difference 2.5\,\si{\degreeCelsius} (45.3\,\si{\degreeCelsius} over the bare ropes against 42.9\,\si{\degreeCelsius} under the greenery), west-oriented pergolas, Mediterranean climate. \\
 & \texttt{ouyang2024} &
Field measurement of shading structures on both axes: air temperature reduced 1.42\,\si{\degreeCelsius} (natural) and 1.31\,\si{\degreeCelsius} (artificial), while mean radiant temperature was reduced 15.93\,\si{\degreeCelsius} and 13.71\,\si{\degreeCelsius} respectively --- an order of magnitude apart. \\
 & \texttt{colter2019} &
Sparse and open canopies group with solid constructed shade structures rather than with dense canopy: PET relief under Parkinsonia, Prosopis and ramadas was 2.9--4.3\,\si{\degreeCelsius} worse than under dense-canopy species. \\
\addlinespace
Productive tree canopy & \texttt{kumar2024} &
Review of 51 green-blue infrastructure types across 202 publications: allotments and city farms are among the least-studied categories, with insufficient data to quantify their cooling. \\
 & \texttt{ziter2019} &
Cooling strengthens above approximately 40\% canopy cover; orchard and agroforestry canopies are deliberately kept open for access and yield, so they sit below that inflection. \\
 & \texttt{colter2019} &
Sparse canopies perform no better than solid artificial shade, which bounds an open production canopy below a closed woodland canopy. \\
 & \texttt{cheung2022} &
Irrigation sustains rather than raises performance: measured agricultural cooling was 0.09--0.43\,\si{\degreeCelsius} against modelled potentials of 2.1--2.5\,\si{\degreeCelsius}, and irrigation increases the daily minimum air temperature as stored heat is released at night. \\
\addlinespace
Productive ground-level planting & \texttt{kumar2024} &
Allotments and city farms are among the least-studied of 51 green-blue infrastructure types, with insufficient data to quantify their cooling. \\
 & \texttt{rost2020} &
The best available allotment-garden measurement: 13 Berlin complexes averaged 2.7\,\si{\kelvin} cooler than built-up areas AT NIGHT. Nocturnal only by design, so it cannot support a daytime value. \\
 & \texttt{armson2012} &
Ground-level vegetation cuts surface temperature substantially while barely affecting globe temperature, which is the basis of the inherited envelope. \\
 & \texttt{cheung2022} &
Measured agricultural cooling 0.09--0.43\,\si{\degreeCelsius} against modelled potentials of 2.1--2.5\,\si{\degreeCelsius}; irrigation raises the daily minimum air temperature at night. \\
\addlinespace
\end{longtable}

\subsection{Calibration notes recorded in configuration}

\begin{longtable}{@{}>{\raggedright\arraybackslash}p{30mm}
                    >{\raggedright\arraybackslash}p{95mm}@{}}
\caption[Archetype calibration notes]{Calibration notes carried in the library,
explaining why an adopted value departs from a naive reading of its sources.
The dense tree canopy, blue-green corridor, and riparian restoration carry no
note; their values follow directly from converging sources.}
\label{tab:app-notes}\\
\toprule
\textbf{Archetype} & \textbf{Note} \\
\midrule
\endfirsthead
\toprule
\textbf{Archetype} & \textbf{Note} \\
\midrule
\endhead
\midrule
\multicolumn{2}{r@{}}{\small\itshape continued on next page}\\
\endfoot
\bottomrule
\endlastfoot
Street tree canopy &
Upper bound capped at 3.0\,\si{\degreeCelsius} rather than the tropical central estimate of 4.0\,\si{\degreeCelsius}, which carries a $\pm$1.6\,\si{\degreeCelsius} interval and reflects favourable conditions. \\
\addlinespace
Continuous shaded corridor &
Scored above street trees because shade is placed where exposure occurs. The narrow-street effect reported by \texttt{keravec2026} is an urban-form effect, not an NbS intervention, and is deliberately not inherited. \\
\addlinespace
Small green area with trees &
Explicitly the small end of the park distribution, bounded below the general park estimate by the size relationship reported in \texttt{bowler2010}. That relationship is qualitative and hedged in the source, and is worded here to match (corrected at 2026.08.04). \\
\addlinespace
Established park &
Represents added cooling over an already partly green baseline, so the achievable increment is bounded below new-forest values. \\
\addlinespace
Extensive green roof &
Benefits are concentrated at roof level and inside the building beneath. The tool states in its output that green-roof benefits are principally building-level rather than street-level. \\
\addlinespace
Green facade or living wall &
Sources conflict outright: the synthesis ranks green walls near the top on air temperature while the micro-scale comparison ranks them below trees on comfort, plausibly explained by near-wall measurement position. Low confidence, conservative envelope, explicit caveat in tool output. \texttt{kumar2024} reports green walls at 4.1 $\pm$ 4.2\,\si{\degreeCelsius}, which is deliberately not adopted: its own standard deviation exceeds its mean, and adopting it selectively would rate a wall above measured tree canopy. The same review's vegetated-balcony figure of 3.8 $\pm$ 2.7\,\si{\degreeCelsius} is declined on the same grounds, which is why the vertical-forest building added at \methver{} inherits this class rather than a canopy class. Adopting the review wholesale would be a full library recalibration with its own sensitivity analysis, and is recorded as a candidate for a future methodology review rather than done piecemeal here. \\
\addlinespace
Bioretention and drainage planting &
No source quantifies air-temperature cooling for rain gardens or bioswales specifically. Value is a reasoned lower bound from adjacent evidence. Primary benefits are stormwater and biodiversity, not cooling. \\
\addlinespace
Blue-green corridor &
The corridor combines an open water body with continuous canopy, which is why it sits above the water-only archetype: its envelope rests on the canopy evidence as much as on the water evidence. \\
\addlinespace
Riparian restoration &
As with the blue-green corridor, the envelope rests on the combination of flowing water and riparian canopy, not on the water surface alone. \\
\addlinespace
Permeable shaded hardscape &
The permeable-pavement estimate rests on a narrower evidence base than the vegetation figures, so the adopted upper bound is set below it. The \texttt{keravec2026} shading-devices figure is admissible HERE, at plaza scale, because that category aggregates area-scale shading; it is deliberately not inherited by the pergola-scale archetype below. \\
\addlinespace
Enclosed courtyard greening &
Small, enclosed, often sheltered from ventilation. Bounded by scale limitation; no courtyard-specific synthesis exists. \\
\addlinespace
Large open water body &
The envelope floor is 0.1\,\si{\degreeCelsius}, not the 0.5\,\si{\degreeCelsius} proposed in the v1.2 review pack. The pack attributed a cooling magnitude near 2\,\si{\degreeCelsius} to \texttt{yao2023a}; retrieval at implementation established that the paper reports no such figure and instead measures 0.1--0.6\,\si{\degreeCelsius} by day. Holding the floor at 0.5 would have made the tool claim more than its most direct field evidence supports, which is exactly the inflation the conservative-calibration rule exists to prevent: where sources conflict, the tool adopts the range its weakest-claiming direct measurement supports and states the conflict rather than smoothing it. The ceiling stays at 3.0\,\si{\degreeCelsius} on \texttt{keravec2026} and the \texttt{volker2013} confidence interval. The width of the envelope is real and is the finding: the remote-sensing-inclusive meta-analysis and the direct field measurements disagree by an order of magnitude, and \texttt{ampatzidis2020} names remote sensing as the overestimating mode. The archetype therefore carries an output caveat stating the disagreement, rather than narrowing the envelope by picking a side the evidence does not support picking. \\
\addlinespace
Small constructed water feature &
THIS ARCHETYPE MUST NEVER BE INTERPOLATED WITH \texttt{large\_water\_body} BY AREA. The two are separate bodies of study, not two points on a dose--response curve: the size-threshold literature is land-surface-temperature only, and no published curve connects them on air temperature. Without this split, constructed wetlands, retention ponds, detention basins and water squares would have inherited \texttt{blue\_green\_corridor} at 1.0--3.0\,\si{\degreeCelsius} on family resemblance and been overstated roughly threefold. The floor is a genuine zero because \texttt{jacobs2020} concludes the daytime effect is negligible in design practice. \\
\addlinespace
Non-canopy vegetation &
\texttt{armson2012} is decisive because it measures both metrics on the same plots and so isolates the surface/air conflation that inflates most published grass figures: cutting surface temperature by 24\,\si{\degreeCelsius} while barely moving globe temperature is the whole finding. The floor is a genuine zero for dry or dormant planting, on \texttt{gill2013} and \texttt{kraemer2022}. \\
\addlinespace
Vegetated shade structure &
THE \texttt{keravec2026} "SHADING DEVICES 2.0\,\si{\degreeCelsius}" FIGURE IS DELIBERATELY NOT USED HERE. That category aggregates area-scale shading, where the cooled air volume persists against advection; a pergola shades roughly \qty{10}{\metre\squared}. Applying it would rate a pergola above measured tree canopy, which \texttt{ouyang2024} measures at 1.42\,\si{\degreeCelsius}. The same figure IS used by the \texttt{permeable\_shaded\_hardscape} archetype, where the scale matches. Low confidence, with $n=1$ for the actual intervention. The ceiling of 2.5\,\si{\degreeCelsius} is \texttt{chafer2020}'s measured peak; note that paper's abstract headlines "up to 5\,\si{\degreeCelsius}", which is greenery against UNSHADED conditions rather than against the bare pergola, and is not the like-for-like contrast this envelope needs. \\
\addlinespace
Productive tree canopy &
Derived, not retrieved: bounded below \texttt{dense\_tree\_canopy} on the \texttt{ziter2019} canopy-cover inflection and the \texttt{colter2019} sparse-canopy finding, and given the same envelope as \texttt{established\_park} with a lower evidence rating. Irrigation is NOT credited as extra degrees --- the \texttt{cheung2022} contrast is specific to agricultural fields, and the same paper measures \SIrange{-2}{-1}{\degreeCelsius} in irrigated urban parks, so no general claim that models overstate irrigation cooling is made here. \\
\addlinespace
Productive ground-level planting &
Derived: inherits the \texttt{non\_canopy\_vegetation} envelope, because cultivated ground-level planting is ground-level vegetation whose cover is periodically broken between crops. Rated low rather than medium because the inheritance is an argument rather than a measurement of this class. The evidence gap is CITABLE (\texttt{kumar2024}) rather than merely unfilled, which is a stronger position than the rain-garden precedent where no source was found at all. \\
\addlinespace
\end{longtable}

\subsection{The \texorpdfstring{\num{121}}{121} catalogue entries and the
archetype each inherits}

Entries are grouped by the source catalogue's own families, with its
element/composite classification reproduced. Neither the family nor the
classification enters a formula. The identifier is the entry's identifier in
the source catalogue; identifiers prefixed \texttt{S} come from the
supplementary proposal curated at version \methver{}
(\Cref{sec:supplement}), which renumbers from scratch and would otherwise
collide with the original catalogue on seventeen identifiers.

{\small
\begin{longtable}{@{}l>{\raggedright\arraybackslash}p{52mm}
                    >{\raggedright\arraybackslash}p{42mm}l@{}}
\caption[The 121 catalogue entries]{The \num{121} curated catalogue entries,
the cooling archetype each inherits, and the source catalogue's
element/composite classification. An entry carries identity, family,
availability, and its primary cooling mechanism; every performance value comes
from the archetype named here.}
\label{tab:app-entries}\\
\toprule
\textbf{ID} & \textbf{Entry} & \textbf{Archetype inherited} & \textbf{Kind} \\
\midrule
\endfirsthead
\multicolumn{4}{@{}l}{\small\itshape \Cref{tab:app-entries} continued}\\
\toprule
\textbf{ID} & \textbf{Entry} & \textbf{Archetype inherited} & \textbf{Kind} \\
\midrule
\endhead
\midrule
\multicolumn{4}{r@{}}{\small\itshape continued on next page}\\
\endfoot
\bottomrule
\endlastfoot
\multicolumn{4}{@{}l}{\textit{Tree-based} (13)}\\
A1 & Strategic individual-tree planting & Street tree canopy & element \\
A2 & Tree cluster & Small green area with trees & element \\
A3 & Tree grove & Dense tree canopy & element \\
A4 & Conventional individual tree pit & Street tree canopy & element \\
A5 & Enlarged tree pit & Street tree canopy & element \\
A6 & Structural-soil tree pit & Street tree canopy & element \\
A7 & Suspended-pavement or soil-cell tree pit & Street tree canopy & element \\
A8 & Stormwater tree pit & Bioretention and drainage planting & element \\
A9 & Bioretention tree pit & Bioretention and drainage planting & element \\
A11 & Continuous tree trench & Street tree canopy & element \\
A12 & Continuous open planting bed for trees & Street tree canopy & element \\
A13 & Connected multi-layer planting bed & Street tree canopy & element \\
A14 & Tree island & Small green area with trees & element \\
\addlinespace
\multicolumn{4}{@{}l}{\textit{Street} (7)}\\
2.2 & Tree Avenue & Street tree canopy & composite \\
2.5 & Green Street & Street tree canopy & composite \\
2.9 & Green Median & Non-canopy vegetation & composite \\
2.13 & Green Tram Corridor & Street tree canopy & composite \\
2.15 & Green Pedestrian Street & Continuous shaded corridor & composite \\
2.16 & Green Alley & Street tree canopy & composite \\
2.17 & Green Parking Lane & Street tree canopy & composite \\
\addlinespace
\multicolumn{4}{@{}l}{\textit{Park} (9)}\\
4.1 & Pocket Park & Established park & composite \\
4.2 & Neighbourhood Park & Established park & composite \\
4.4 & District Park & Established park & composite \\
4.5 & Regional Park & Established park & composite \\
4.11 & Water Park & Established park & composite \\
4.15 & Naturalistic Park & Established park & composite \\
4.17 & Stormwater Park & Established park & composite \\
4.20 & Linear Park & Established park & composite \\
4.22 & Coastal Park & Established park & composite \\
\addlinespace
\multicolumn{4}{@{}l}{\textit{Public space} (6)}\\
3.1 & Tree Plaza & Small green area with trees & composite \\
3.2 & Cooling Plaza & Permeable shaded hardscape & composite \\
3.3 & Permeable Plaza & Permeable shaded hardscape & composite \\
3.4 & Civic Green & Established park & composite \\
3.5 & Courtyard Greening & Enclosed courtyard greening & composite \\
3.15 & Green Playground & Small green area with trees & composite \\
\addlinespace
\multicolumn{4}{@{}l}{\textit{Woodland and forest} (6)}\\
5.1 & Urban woodland (site) & Dense tree canopy & composite \\
5.3 & Microforest & Dense tree canopy & composite \\
5.10 & Degraded Woodland Restoration & Dense tree canopy & composite \\
5.15 & Afforestation & Dense tree canopy & composite \\
5.16 & Reforestation & Dense tree canopy & composite \\
5.18 & Urban Woodland Buffer & Dense tree canopy & composite \\
\addlinespace
\multicolumn{4}{@{}l}{\textit{Ground level} (11)}\\
B1 & Conventional planting bed & Non-canopy vegetation & element \\
B2 & Shrub-dominated planting bed & Non-canopy vegetation & element \\
B3 & Herbaceous perennial bed & Non-canopy vegetation & element \\
B4 & Groundcover planting & Non-canopy vegetation & element \\
B6 & Pocket meadow & Non-canopy vegetation & element \\
B7 & Rain garden & Bioretention and drainage planting & element \\
B8 & Bioswale & Bioretention and drainage planting & element \\
B9 & Infiltration planting bed & Bioretention and drainage planting & element \\
B10 & Depressed bioretention bed & Bioretention and drainage planting & element \\
B11 & Vegetated planter with stormwater function & Bioretention and drainage planting & element \\
B12 & Large-volume connected planters & Bioretention and drainage planting & element \\
\addlinespace
\multicolumn{4}{@{}l}{\textit{Water landscape} (16)}\\
6.1 & Constructed Wetland & Small constructed water feature & composite \\
6.2 & Retention Pond & Small constructed water feature & composite \\
6.3 & Detention Basin & Small constructed water feature & composite \\
6.4 & Water Square & Small constructed water feature & composite \\
6.5 & Daylighted Stream & Riparian restoration & composite \\
6.7 & River Restoration & Riparian restoration & composite \\
6.8 & Canal Restoration & Riparian restoration & composite \\
6.10 & Floodplain Restoration & Riparian restoration & composite \\
6.11 & Urban Lake Restoration & Large open water body & composite \\
6.12 & Living Shoreline & Large open water body & composite \\
6.13 & Floating Wetland & Small constructed water feature & element \\
6.14 & Riparian Buffer & Riparian restoration & element \\
S12A.1 & Subsurface-flow treatment wetland & Non-canopy vegetation & composite \\
S17.4 & Mangrove Restoration & Riparian restoration & composite \\
S17.6 & Dune Restoration & Non-canopy vegetation & composite \\
S17.7 & Coastal Wetland Restoration & Non-canopy vegetation & composite \\
\addlinespace
\multicolumn{4}{@{}l}{\textit{Ecological network} (6)}\\
7.1 & Green Corridor & Blue-green corridor & composite \\
7.2 & Ecological Corridor & Blue-green corridor & composite \\
7.4 & Blue-Green Corridor & Blue-green corridor & composite \\
7.5 & River Corridor & Riparian restoration & composite \\
7.6 & Railway Green Corridor & Street tree canopy & composite \\
7.14 & Coastal Ecological Corridor & Large open water body & composite \\
\addlinespace
\multicolumn{4}{@{}l}{\textit{District} (5)}\\
10.1 & Green Infrastructure Network & Blue-green corridor & composite \\
10.2 & Blue-Green Infrastructure Network & Blue-green corridor & composite \\
10.3 & Urban Cooling Network & Dense tree canopy & composite \\
10.4 & Green Mobility Network & Street tree canopy & composite \\
10.10 & Forest District & Dense tree canopy & composite \\
\addlinespace
\multicolumn{4}{@{}l}{\textit{Green roof} (12)}\\
C1 & Extensive green roof & Extensive green roof & element \\
C2 & Semi-intensive green roof & Extensive green roof & element \\
C3 & Intensive green roof & Small green area with trees & element \\
C5 & Biodiverse green roof & Extensive green roof & element \\
C6 & Brown roof & Extensive green roof & element \\
C8 & Blue-green roof & Extensive green roof & element \\
C11 & Biosolar roof & Extensive green roof & element \\
C13 & Productive green roof & Productive ground-level planting & element \\
C14 & Rooftop farm & Productive ground-level planting & element \\
C15 & Roof meadow & Non-canopy vegetation & element \\
C16 & Roof wetland & Small constructed water feature & element \\
C17 & Tree roof or rooftop woodland & Dense tree canopy & element \\
\addlinespace
\multicolumn{4}{@{}l}{\textit{Vertical greening} (10)}\\
D1 & Direct ground-rooted green façade & Green facade or living wall & element \\
D2 & Indirect ground-rooted green façade & Green facade or living wall & element \\
D5 & Planter-rooted green façade & Green facade or living wall & element \\
D6 & Cascading façade planting & Green facade or living wall & element \\
D7 & Modular living wall & Green facade or living wall & element \\
D8 & Substrate-based living wall & Green facade or living wall & element \\
D9 & Hydroponic living wall & Green facade or living wall & element \\
D11 & Vegetated screen & Green facade or living wall & element \\
D12 & Green noise barrier & Green facade or living wall & element \\
D13 & Vegetated retaining wall & Green facade or living wall & element \\
\addlinespace
\multicolumn{4}{@{}l}{\textit{Building} (5)}\\
9.29 & Green Balcony & Green facade or living wall & element \\
9.30 & Green Terrace & Extensive green roof & element \\
9.33 & Green Podium & Small green area with trees & composite \\
S9.19 & Green Loggia & Enclosed courtyard greening & element \\
S9.22 & Vertical Forest Building & Green facade or living wall & composite \\
\addlinespace
\multicolumn{4}{@{}l}{\textit{Vegetated shelter} (3)}\\
E1 & Green pergola & Vegetated shade structure & element \\
E3 & Vegetated arcade & Continuous shaded corridor & element \\
E6 & Green waiting shelter & Vegetated shade structure & element \\
\addlinespace
\multicolumn{4}{@{}l}{\textit{Productive landscape} (7)}\\
8.1 & Community Garden & Productive ground-level planting & composite \\
8.2 & Allotment Garden & Productive ground-level planting & composite \\
8.3 & Urban Orchard & Productive tree canopy & composite \\
8.4 & Food Forest & Productive tree canopy & composite \\
8.5 & Urban Farm & Productive ground-level planting & composite \\
8.9 & School Garden & Productive ground-level planting & composite \\
8.15 & Agroforestry System & Productive tree canopy & composite \\
\addlinespace
\multicolumn{4}{@{}l}{\textit{Soil, slope and site restoration} (5)}\\
SA16 & Depaving & Non-canopy vegetation & element \\
S11B.16 & Vegetated Terracing & Non-canopy vegetation & composite \\
S11B.17 & Natural Slope Restoration & Non-canopy vegetation & composite \\
S19.4 & Phytoremediation & Non-canopy vegetation & composite \\
S21.8 & Brownfield Ecological Succession & Non-canopy vegetation & composite \\
\addlinespace
\end{longtable}
}

\subsection{The twenty-one suitability overrides}

An entry overrides an inherited suitability condition only where the source
catalogue's own description of the entry makes the inherited value plainly
wrong (\Cref{sec:suitability}). No new number enters the methodology through an
override: every replacement value is one already present in the cited library.
A test fails if this list grows past \num{25} entries.

\begin{longtable}{@{}>{\raggedright\arraybackslash}p{52mm}
                    >{\raggedright\arraybackslash}p{26mm}
                    >{\raggedright\arraybackslash}p{50mm}@{}}
\caption[Per-entry suitability overrides]{The \num{21} per-entry suitability
overrides, with the reason the inherited value is wrong.}
\label{tab:app-overrides}\\
\toprule
\textbf{Entries} & \textbf{Override} & \textbf{Why the inherited value is
wrong} \\
\midrule
\endfirsthead
\toprule
\textbf{Entries} & \textbf{Override} & \textbf{Why the inherited value is
wrong} \\
\midrule
\endhead
\bottomrule
\endlastfoot
Tree grove, microforest, rooftop woodland &
area $\rightarrow$ \qty{200}{\metre\squared} &
The dense tree canopy's \qty{2000}{\metre\squared} describes the woodland its
\emph{cooling evidence} comes from, not the footprint of a site element. The
catalogue defines a microforest as a small dense patch on a constrained
site. \\
\addlinespace
Intensive green roof, productive green roof, roof meadow, roof wetland
(\qty{50}{\metre\squared}); rooftop farm, rooftop woodland
(\qty{200}{\metre\squared}) &
area, and soil $\rightarrow$ none &
A roof's growing medium is \emph{built}, not ground. Requiring ground soil on a
roof would flag every roof entry whose archetype happens to be a ground-level
class. \\
\addlinespace
Green podium, water square, floating wetland &
soil $\rightarrow$ none &
Built structure or open water, likewise. The water square is additionally set
to \qty{200}{\metre\squared}, at plaza scale. \\
\addlinespace
Direct and indirect ground-rooted green fa\c{c}ades, green noise barrier,
vegetated retaining wall &
soil $\rightarrow$ limited &
The catalogue defines these as rooted in \emph{natural ground}, which the
containerised living-wall archetype does not require. This override makes the
tool \emph{stricter}, not more permissive. \\
\addlinespace
Pocket park &
area $\rightarrow$ \qty{200}{\metre\squared} &
This entry \emph{is} the tool's own cited pocket park; its minimum area is the
one the previous library established. \\
\addlinespace
Riparian buffer &
area $\rightarrow$ \qty{200}{\metre\squared} &
The catalogue distinguishes a discrete vegetated strip from the restored
corridor the archetype describes. \\
\addlinespace
Railway green corridor, green mobility network &
area $\rightarrow$ \qty{1000}{\metre\squared} &
Corridor- and district-scale systems inheriting an element archetype. \\
\addlinespace
Urban farm (\qty{500}{\metre\squared}), agroforestry system
(\qty{1000}{\metre\squared}) &
area &
The catalogue distinguishes field-scale production from garden scale. \\
\end{longtable}

\clearpage
\section{Complete weight tables}
\label{app:weights}

All weights are from \texttt{config/weights.yaml} at version \methver{}. Within
each block the weights sum to unity. All are expert judgment in the sense of
\textcite{nardo2008}; where a source informs the \emph{selection} of indicators
or the \emph{relative ordering} of weights, it is named.

\subsection{Heat Exposure Score (data-rich path only)}

\begin{center}\small
\begin{tabular}{@{}lr@{}}
\toprule
\textbf{Component} & \textbf{Weight} \\
\midrule
Land surface temperature anomaly & 0.40 \\
Imperviousness                   & 0.25 \\
Solar exposure                   & 0.20 \\
Vegetation deficit               & 0.15 \\
\midrule
\textbf{Total} & \textbf{1.00} \\
\bottomrule
\end{tabular}
\end{center}

\noindent\emph{Basis.} Component selection reflects established
surface-energy-balance drivers of urban heat. \textcite{ziter2019} gives direct
empirical support for two components: impervious cover raises air temperature
approximately linearly and canopy cover reduces it nonlinearly. The relative
weights are expert judgment.

\subsection{Vulnerability Score}

\begin{center}\small
\begin{tabular}{@{}lr@{}}
\toprule
\textbf{Component} & \textbf{Weight} \\
\midrule
Population density        & 0.40 \\
Vulnerable groups         & 0.40 \\
Cooling access deficit    & 0.20 \\
\midrule
\textbf{Total} & \textbf{1.00} \\
\bottomrule
\end{tabular}
\end{center}

\noindent\emph{Basis.} Indicator selection follows \textcite{reid2009}, whose
factor analysis of ten heat-vulnerability variables found four factors
explaining over \SI{75}{\percent} of variance, with air-conditioning prevalence
emerging as an independent dimension. Density and vulnerable-group presence are
weighted equally because that analysis gives no basis for ranking one above the
other; cooling access carries less weight as a single, frequently unknown
indicator.

\subsection{Heat Priority Index and site adjustment}

\begin{center}\small
\begin{tabular}{@{}lr@{}lr@{}}
\toprule
\multicolumn{2}{@{}l}{\textbf{Heat Priority Index}} &
\multicolumn{2}{l}{\textbf{Site adjustment}} \\
\midrule
Heat exposure  & 0.60 & \quad Canopy      & 0.40 \\
Vulnerability  & 0.40 & \quad Soil--water & 0.25 \\
               &      & \quad Scale       & 0.20 \\
               &      & \quad Climate     & 0.15 \\
\midrule
\textbf{Total} & \textbf{1.00} & \quad\textbf{Total} & \textbf{1.00} \\
\bottomrule
\end{tabular}
\end{center}

\noindent\emph{Basis for site adjustment.} Canopy carries the largest weight as
the condition with the strongest direct empirical support
\parencite{ziter2019}. The climate component is justified by
\textcite{keravec2026}: adding K\"oppen--Geiger subzone raises variance
explained in cooling effectiveness from \SI{12}{\percent} to
\SI{78}{\percent}.

\subsection{Composite sub-scores}

\begin{center}\small
\begin{tabular}{@{}lr@{\qquad}lr@{\qquad}lr@{}}
\toprule
\multicolumn{2}{@{}l}{\textbf{NbS Suitability}} &
\multicolumn{2}{l}{\textbf{Co-benefits}} &
\multicolumn{2}{l}{\textbf{Equity}} \\
\midrule
Space         & 0.30 & Biodiversity     & 0.25 & Vulnerable user benefit & 0.35 \\
Soil          & 0.25 & Stormwater       & 0.25 & Public accessibility    & 0.25 \\
Water         & 0.20 & Public health    & 0.20 & Safety and comfort      & 0.20 \\
Maintenance   & 0.15 & Social inclusion & 0.15 & Participation relevance & 0.20 \\
Urban context & 0.10 & Urban quality    & 0.15 &                         &      \\
\midrule
\textbf{Total} & \textbf{1.00} & \textbf{Total} & \textbf{1.00} &
\textbf{Total} & \textbf{1.00} \\
\bottomrule
\end{tabular}
\end{center}

\noindent The rules mapping user inputs to these individual sub-indicator
scores were fixed at version 2026.08.01 in
\texttt{config/derived\_scores.yaml} and reproduced in
\Cref{tab:suitability-rules}. The Equity Score does not enter the final
aggregation (\Cref{sec:composite-subscores}). The derived cost feasibility
score combines the payback bracket with two delivery-risk modifiers
(\Cref{eq:cost-feasibility-rule}):

\begin{center}\small
\begin{tabular}{@{}lr@{}}
\toprule
\textbf{Cost feasibility component} & \textbf{Weight} \\
\midrule
Payback bracket score            & 0.50 \\
Implementation complexity (inverted) & 0.25 \\
Maintenance intensity (inverted)     & 0.25 \\
\midrule
\textbf{Total} & \textbf{1.00} \\
\bottomrule
\end{tabular}
\end{center}

\subsection{Final aggregation and effective weights}

\begin{center}\small
\begin{tabular}{@{}lrr@{}}
\toprule
\textbf{Component} & \textbf{Nominal} & \textbf{Effective} \\
\midrule
Heat Priority Index    & 0.25 & --- \\
\quad heat exposure    &      & 0.15 \\
\quad vulnerability    &      & 0.10 \\
Cooling Potential      & 0.25 & 0.25 \\
NbS Suitability        & 0.15 & 0.15 \\
Vulnerability (direct) & 0.15 & 0.15 \\
Co-benefits            & 0.10 & 0.10 \\
Cost Feasibility       & 0.10 & 0.10 \\
\midrule
\textbf{Total}                & \textbf{1.00} & \textbf{1.00} \\
\midrule
\multicolumn{3}{@{}l}{\itshape Restated, not added:}\\
\textbf{Vulnerability, total} &      & \textbf{0.25} \\
\textbf{Heat exposure, total} &      & \textbf{0.15} \\
\bottomrule
\end{tabular}
\end{center}

\noindent The equity-forward double contribution of vulnerability is deliberate
and is argued in \Cref{sec:effective-weights}. The sensitivity-analysis
perturbation parameter is $\pm\,\SI{25}{\percent}$
(\Cref{sec:sensitivity}).

\clearpage
\section{Input field reference}
\label{app:fields}

\begin{keynote}
\noindent The definitive enumeration is fixed by the implemented input schema
(\texttt{nature\_cooling.\allowbreak engine.\allowbreak models.\allowbreak
 AssessmentInput}). Fields marked
\emph{Required} are exactly those without which a documented formula cannot be
evaluated at all; every other field is \emph{Optional}, and its absence is
handled through a disclosed default or an explicit \emph{not calculated}
status, itemised in the result.
\end{keynote}

\begin{longtable}{@{}>{\raggedright\arraybackslash}p{46mm}
                    >{\raggedright\arraybackslash}p{22mm}
                    >{\raggedright\arraybackslash}p{60mm}@{}}
\caption[Input field reference]{Input fields, their status, and the
normalisation applied to each.}
\label{tab:app-fields}\\
\toprule
\textbf{Field} & \textbf{Status} & \textbf{Normalisation / use} \\
\midrule
\endfirsthead
\multicolumn{3}{@{}l}{\small\itshape \Cref{tab:app-fields} continued}\\
\toprule
\textbf{Field} & \textbf{Status} & \textbf{Normalisation / use} \\
\midrule
\endhead
\midrule
\multicolumn{3}{r@{}}{\small\itshape continued on next page}\\
\endfoot
\bottomrule
\endlastfoot

\multicolumn{3}{@{}l}{\textit{Site characteristics}}\\
\fieldname{site_area_m2} & Required &
Denominator of the canopy condition and of shade potential
(\Cref{eq:shade}). \\
\fieldname{existing_tree_canopy_percent} & Optional &
Canopy condition derivation (\Cref{tab:derivation}). \\
\fieldname{existing_green_cover_percent} & Optional &
Vegetation deficit, $\clamp{100-g}$ (\Cref{eq:vegdef}). \\
\fieldname{impervious_surface_percent} & Optional &
Percentage used directly as a \numrange{0}{100} score in
\Cref{eq:heat-exposure-rich}; the identity mapping is declared in
configuration. \\
\fieldname{soil_availability} & Optional &
none 25 / limited 50 / moderate 75 / high 100; unknown 50. Also feeds the
soil--water condition. \\
\fieldname{irrigation_availability} & Optional &
none 25 / occasional 50 / reliable 100; unknown 50. Also feeds the soil--water
condition. \\
\fieldname{current_shade_level} & Optional &
very low 25 / low 50 / medium 75 / high 100; unknown 50. Counts toward
cooling-block completeness; enters no formula at version \methver{}. \\
\fieldname{land_use} & Optional &
Urban-context suitability sub-indicator: in the entry's own land-use list 100,
outside it 25, unknown 50 (\Cref{tab:suitability-rules}). That list is the
availability matrix's list, so an entry is in context exactly where it is
offered (\Cref{sec:availability}). \emph{campus} and \emph{memorial} were
added at version 2026.08.03. \\
\addlinespace

\multicolumn{3}{@{}l}{\textit{Availability conditions --- these gate which
interventions are offered and feed \textbf{no} score}}\\
\fieldname{includes_railway} & Optional &
yes / no. Gates exactly one entry, the railway green corridor. Unanswered
withholds it. \textbf{Availability only}: enters no formula, no confidence
block, and no aggregation (\Cref{sec:availability}). New at version
\methver{}. \\
\fieldname{existing_woodland} & Optional &
yes / no. Gates woodland \emph{restoration} entries only; woodland
\emph{creation} entries are ungated. Unanswered withholds the restoration
entries. \textbf{Availability only}. New at version \methver{}. \\
\fieldname{waterfront_type} & Optional &
lake / river / dry river / covered river / sea. Gates twelve water-landscape
and ecological-network entries by category; the four \emph{constructed} water
features carry no waterfront condition. Unanswered withholds the gated
entries. \textbf{Availability only}. New at version \methver{}. \\
\fieldname{productive_governance} & Optional &
Multi-select over community / individual / institutional / commercial. Filters
the seven productive entries by who delivers them; an \emph{empty or
unanswered} selection suppresses nothing, because absence of an answer is not
evidence that no delivery model exists. Canonicalised on input so that two
selections of the same models produce byte-identical results.
\textbf{Availability only}. New at version \methver{}. \\
\addlinespace

\multicolumn{3}{@{}l}{\textit{Project information}}\\
\fieldname{assessment_scale} & Required &
Scale condition: city, district $\rightarrow$ excellent; neighbourhood
$\rightarrow$ good; site, building $\rightarrow$ moderate. \\
\fieldname{country} & Optional &
Grid emission factor lookup in \texttt{country\_defaults.yaml}
(\Cref{sec:ghg}). \\
\addlinespace

\multicolumn{3}{@{}l}{\textit{Climate and heat exposure}}\\
\fieldname{climate_zone} & Required &
Climate condition lookup (\Cref{tab:climate-matrix}). \\
\fieldname{lst_anomaly_c} & Optional &
Linear, \qty{0}{\degreeCelsius}$\rightarrow$0,
\qty{+10}{\degreeCelsius}$\rightarrow$100, clamped (\Cref{eq:lst}). Selects the
data-rich heat exposure path. \\
\fieldname{heat_exposure_level} & Optional &
Standard levels. Used directly as the Heat Exposure Score on the data-poor
path. \\
\fieldname{solar_exposure} & Optional & Standard levels. \\
\fieldname{heat_index_concern} & Optional & Standard levels. \\
\addlinespace

\multicolumn{3}{@{}l}{\textit{Vulnerability and equity}}\\
\fieldname{population_density} & Optional & Standard levels. \\
\fieldname{vulnerable_population_presence} & Optional & Standard levels. \\
\fieldname{access_to_cooled_indoor_space} & Optional &
\textbf{Inverted} levels: high 25 / medium 50 / low 100; unknown 50
(\Cref{sec:inverted}). \\
\fieldname{access_to_cool_outdoor_refuge} & Optional &
\textbf{Inverted} levels, as above. Second indicator of the same cooling
access deficit dimension; the two are averaged over those supplied and share
one confidence slot. New at version 2026.08.03. \\
\fieldname{safety_concern} & Optional &
Standard levels; equity safety-and-comfort sub-indicator (deficit reading). \\
\fieldname{public_accessibility} & Optional &
Standard levels; equity sub-indicator. New at version 2026.08.01. \\
\fieldname{community_participation} & Optional &
Standard levels; equity participation-relevance sub-indicator. New at version
2026.08.01. \\
\addlinespace

\multicolumn{3}{@{}l}{\textit{NbS intervention}}\\
\fieldname{nbs_type} & Required &
A \emph{list} since version \methver{}: one or more catalogue entries.
Selects each entry, the archetype it inherits, and that archetype's
climate-family row. Duplicates are rejected; selection order is preserved in
the itemised component results and changes no output. More than five
components raises a warning, not an error (\Cref{sec:packages}). \\
\fieldname{new_canopy_area_at_maturity_m2} & Optional &
Canopy condition and shade potential. \\
\fieldname{expected_maturity_period_years} & Optional &
Time-to-benefit class (immediate / short / medium / long term). No maturity
discount is applied to cooling values. \\
\fieldname{intervention_area_m2} & Optional &
Validation warning when exceeding the site area; enters no score. \\
\fieldname{implementation_complexity} & Optional &
Inverted levels; cost feasibility (\Cref{eq:cost-feasibility-rule}) and
investment readiness. \\
\fieldname{maintenance_intensity} & Optional &
Inverted levels; cost feasibility and the suitability maintenance
sub-indicator. \\
\addlinespace

\multicolumn{3}{@{}l}{\textit{Cost and energy (all optional)}}\\
\fieldname{nearby_building_cooling_demand_relevant} & Optional &
Precondition for the energy derivation (\Cref{sec:energy}). \\
\fieldname{annual_cooling_energy_demand_kwh} & Optional &
$D$ in \Cref{eq:energy}. Absent: \texttt{missing\_energy\_demand}. \\
\fieldname{energy_price_per_kwh} & Optional &
$p$ in \Cref{eq:annual-savings}. Absent: cost savings not calculated. \\
\fieldname{capital_cost} & Optional &
$C_{\mathrm{capital}}$ in \Cref{eq:payback}. Entered \textbf{once for the
whole package}: it covers every selected component and is not apportioned
(\Cref{tab:package-rules}). Absent: capital cost, payback, cost feasibility,
and investment readiness all \emph{not estimated}. \\
\fieldname{currency} & Optional &
ISO 4217 code, echoed with cost outputs; never used in scoring. \\
\fieldname{grid_emission_factor_kgco2e_per_kwh} & Optional &
$\phi$ in \Cref{eq:ghg}; a user-supplied factor takes precedence over the
country table and is itemised as an applied assumption. \\
\end{longtable}

\clearpage
\section{Glossary}
\label{app:glossary}

\begin{description}[leftmargin=0pt, labelwidth=0pt, labelsep=0pt, style=nextline, itemsep=4pt]

\item[Air temperature] The temperature of the air, here measured or estimated
at pedestrian level. The quantity this methodology reports. Distinct from land
surface temperature and from thermal comfort indices.

\item[Archetype (cooling archetype)] A cited evidence class carrying every
performance value --- envelope, base score, evidence rating, suitability
conditions, co-benefit defaults --- and the citations behind them. There are
\num{18} (\Cref{sec:archetypes}).

\item[Availability] Which catalogue entries the tool \emph{offers} for a
given site, decided by scale, land use, and the four gating questions of
\Cref{sec:availability}. Availability feeds no score.

\item[Base cooling score] A \numrange{0}{100} index expressing an archetype's
relative expected cooling performance before site adjustment. Not a
temperature.

\item[Catalogue entry (typology)] One of the \num{121} curated entries of the
source catalogue. It carries identity, family, availability, and the one
archetype it inherits, and no performance value of its own.

\item[Block confidence] A property of an \emph{assessment}: how completely the
user supplied the inputs a given output block depends on
(\Cref{tab:confidence}). Distinct from evidence confidence.

\item[Boundary-layer urban heat island] The heat island of the air mass above
roof level, distinguished from the canopy-layer and surface heat islands by
\textcite{oke1976}.

\item[Canopy-layer urban heat island (CUHI)] The heat island of the air between
the ground and roof level --- the layer people occupy. The quantity relevant to
this methodology.

\item[Clamp] $\clamp{x}$ bounds a value to the \numrange{0}{100} reporting
scale.

\item[Clip] $\clip{x}{l}{u}$ bounds a value to a stated interval. Used to
confine reported temperature estimates to the literature envelope
(\Cref{sec:cooling}).

\item[Composite indicator] An index formed by normalising and aggregating
several sub-indicators. Construction here follows \textcite{nardo2008}.

\item[Cooling Potential Score] The Layer 2 headline score: the typology's base
cooling score scaled by the site adjustment factor (\Cref{eq:cps}).

\item[Envelope] The literature-supported interval of temperature reduction for
an archetype. Its lower bound represents a poorly performing instance, its
upper bound a well-executed instance under favourable conditions. Not a
prediction and not a confidence interval.

\item[Evidence confidence] A property of an \emph{archetype}: how well
grounded its adopted values are (high, medium, low). Inherited by every entry
resolving to it, and propagates into an assessment by capping cooling-block
confidence (\Cref{sec:evidence-propagation}).

\item[Golden scenario] A stored input case with a hand-verified expected
output, used as regression armour against unintended methodology drift and as
the basis of the sensitivity analysis.

\item[Heat Priority Index (HPI)] The Layer 1 headline score, combining heat
exposure and vulnerability (\Cref{eq:hpi}). Independent of the proposed
intervention.

\item[K\"oppen--Geiger classification] A climate classification system whose
subzones (Af, BWh, Csa, Cfa, Cfb, Cwa, \dots) raise the variance explained in
cooling effectiveness from \SI{12}{\percent} to \SI{78}{\percent}
\parencite{keravec2026}.

\item[Land surface temperature (LST)] The radiometric temperature of the
ground and building surfaces, typically satellite-derived. Substantially
exceeds air temperature \parencite{venter2021}; used here only as a relative
screening proxy (\Cref{sec:lst}).

\item[NbS Cooling Opportunity Score] The final aggregated score
(\Cref{eq:aggregation}), on a \numrange{0}{100} scale, for comparing options
and sites.

\item[Nature-based solutions (NbS)] Interventions using vegetation, water, and
natural processes to deliver benefits --- here, urban cooling.

\item[Package] An assessment proposing more than one catalogue entry. Each
component is scored and reported individually; the headline outputs are
combined under \Cref{tab:package-rules}, which caps temperature at the
best-evidenced component and never sums it.

\item[PET (physiologically equivalent temperature)] A thermal comfort index
representing perceived thermal conditions. \textbf{Not} air temperature; used
in this methodology only for relative ranking and caveats.

\item[Screening level] An assessment tier producing indicative, comparable
estimates for triage, positioned between unstructured judgment and detailed
simulation.

\item[Site adjustment factor] The composite multiplier $A \in [0.5, 1.2]$
derived from canopy, soil--water, scale, and climate conditions
(\Cref{eq:adjustment}).

\item[Super-additivity] The hypothesis that combined interventions cool more
than their strongest component. Not adopted, because no retrieved source
quantifies it (\Cref{sec:super-additivity}). A package's advantage is expressed
as co-benefit breadth instead.

\item[Surface urban heat island (SUHI)] The heat island measured in surface
temperatures. Reported at roughly six times the canopy-layer magnitude by
\textcite{venter2021}.

\item[UTCI (universal thermal climate index)] A thermal comfort index. As with
PET, not air temperature, and not reported by this tool.

\end{description}

\clearpage
\section{Source verification status}
\label{app:verification}

Every source enters the methodology's bibliography only after its bibliographic
details have been checked against the publisher record, an institutional
repository, or an indexing service. \Cref{tab:app-verification} reproduces the
verification register from \texttt{docs/methodology/BIBLIOGRAPHY.md} at version
\methver{}.

Three verification levels are used.

\begin{description}[leftmargin=1.5em, style=sameline, font=\normalfont\itshape]
  \item[Full] Bibliographic details \emph{and} the specific quantitative
  finding used by the methodology were both confirmed from a retrieved source.
  \item[Metadata + finding (secondary)] Bibliographic details confirmed from
  the publisher or repository; the quantitative finding confirmed from an
  authoritative secondary description --- a publisher abstract page or an
  indexing service --- rather than from the full text.
  \item[Metadata only] Bibliographic details confirmed; the source is
  referenced qualitatively only, and no numeric value in the tool depends on
  it.
\end{description}

\begin{keynote}
\noindent This register is reproduced rather than summarised because the
distinction it records is material. Several publisher sites --- Elsevier, PNAS,
Science --- block automated retrieval, so a number of findings were confirmed
from open repositories, PubMed, or publisher abstract pages. Where a value in
the tool depends on a finding that is not \emph{Full}, the configuration marks
its confidence accordingly.

A reader should not read \emph{metadata + finding (secondary)} as an implied
full-text reading. It means what it says: the number was confirmed from an
authoritative description of the paper, not from the paper.
\end{keynote}

\begin{longtable}{@{}l>{\raggedright\arraybackslash}p{34mm}
                    >{\raggedright\arraybackslash}p{66mm}@{}}
\caption[Source verification register]{Verification status of every source
cited by the methodology.}
\label{tab:app-verification}\\
\toprule
\textbf{Key} & \textbf{Verified} & \textbf{Basis and notes} \\
\midrule
\endfirsthead
\multicolumn{3}{@{}l}{\small\itshape \Cref{tab:app-verification} continued}\\
\toprule
\textbf{Key} & \textbf{Verified} & \textbf{Basis and notes} \\
\midrule
\endhead
\midrule
\multicolumn{3}{r@{}}{\small\itshape continued on next page}\\
\endfoot
\bottomrule
\endlastfoot

\texttt{keravec2026} & Full &
IOPscience, open access. The primary cross-typology anchor and the empirical
basis for treating climate suitability as a first-order adjustment. \\
\addlinespace
\texttt{ziter2019} & Full &
PubMed record, PMID 30910972. \\
\addlinespace
\texttt{zolch2016} & Full &
Lund University research portal record. \emph{Note:} PET is a thermal-comfort
index, not air temperature; used for relative ranking of intervention families
and for the green-roof street-level caveat, never as an air-temperature
value. \\
\addlinespace
\texttt{santamouris2014} & Full &
Open PDF via EEA Climate-ADAPT; abstract read directly. \emph{Critical
caveat:} the \SIrange{0.3}{3}{\kelvin} figure describes city-scale mass
deployment in simulation studies, not one green roof on one building, and is
therefore not used as a site-level typology value. \\
\addlinespace
\texttt{reid2009} & Full &
PMC open access record, PMC2801183. Grounds the vulnerability sub-indicators
and the treatment of low cooling access as an independent vulnerability
dimension; within that dimension, the \emph{indoor} indicator. \\
\addlinespace
\texttt{burkart2016} & Metadata + finding &
PubMed record, PMID 26566198; findings from the \emph{Environmental Health
Perspectives} abstract. Grounds the \emph{outdoor} cooling-refuge indicator
added at version \methver{}: reachable vegetation and water quantified as
modifiers of elderly heat mortality in Lisbon. \\
\addlinespace
\texttt{sera2019} & Metadata + finding &
PubMed record, PMID 30815699; findings from the \emph{International Journal of
Epidemiology} abstract. Corroborates \texttt{burkart2016} across 340 cities in
22 countries, and independently supports population density as a vulnerability
indicator. \\
\addlinespace
\texttt{ember} & Full &
Licence and coverage confirmed on Ember's data pages; 215 countries, CC-BY-4.0,
redistribution permitted with attribution. \\
\addlinespace
\texttt{beck2023} & Full &
Open access at \emph{Scientific Data}; licence, period, resolution and the
30-class legend confirmed against the published archive, whose checksum is
recorded in \texttt{tools/build\_datasets.py}. CC BY 4.0, redistribution
permitted with attribution. \\
\addlinespace
\texttt{naturalearth} & Full &
Terms of use and release v5.1.2 confirmed at source; SHA-256 of both retrieved
scales recorded in \texttt{tools/build\_datasets.py}. Public domain; attribution
recorded although not required. \\
\addlinespace
\texttt{nardo2008} & Full &
Full bibliographic record from the JRC repository; step sequence confirmed from
OECD/JRC descriptions. \\
\addlinespace
\texttt{worldbank2021} & Full &
World Bank plain-text document retrieved directly. The documented
justification for shipping no global default unit costs. \\
\addlinespace
\texttt{jacobs2020} & Full &
Open publisher PDF via the Wageningen repository, read directly. The primary
source for the small constructed water archetype, and the reason constructed
wetlands, ponds, basins, and water squares are not allowed to inherit the
corridor envelope on family resemblance. \\
\addlinespace
\texttt{volker2013} & Full &
Open-access publisher PDF, read directly. \emph{Caveats recorded:} a single
pooled estimate across ponds, lakes, and rivers together, so it cannot justify
a size-dependent gradient; and its magnitude depends on the inclusion of
remote-sensing data, which \texttt{ampatzidis2020} identifies as the
overestimating mode. \\
\addlinespace
\texttt{ampatzidis2020} & Full \emph{(remote-sensing statement: metadata +
finding, secondary)} &
Publisher abstract read in full via the University of Bath research portal for
the diurnal-variability and nocturnal findings; the compact ``remote sensing
overestimates'' phrasing appears in the article highlights via an indexing
service. Governs how the two water archetypes are read. \\
\addlinespace
\texttt{armson2012} & Full &
University of Manchester Research Explorer record, complete abstract. The
decisive source for the non-canopy vegetation archetype, because it measures
surface and globe temperature on the same plots and so isolates the
surface/air conflation. \\
\addlinespace
\texttt{kraemer2022} & Full &
Open access, Frontiers, read directly. \emph{Note:} an earlier draft of the
review pack cited this source in connection with allotment gardens. It is a
study of two Leipzig \emph{parks} under drought and is cited as such
(\Cref{sec:water-body}). \\
\addlinespace
\texttt{chafer2020} & Full &
Publisher PDF retrieved and read; the author list was independently confirmed
against the Crossref record after the review pack recorded it wrongly.
\emph{Caveat recorded:} the paper's abstract headlines ``up to
\qty{5}{\degreeCelsius}'', which is greenery against \emph{unshaded}
conditions; the adopted \qty{2.5}{\degreeCelsius} is the like-for-like contrast
against an identical bare pergola. \\
\addlinespace
\texttt{colter2019} & Full &
Arizona State University Elsevier Pure record, complete abstract.
\emph{Caveat:} PET is a comfort index, used here only for the ordering it
establishes, under narrow conditions --- hot summer midday, desert climate, six
tree taxa. \\
\addlinespace
\texttt{rost2020} & Full &
Publisher-deposited abstract retrieved verbatim. \emph{Critical caveat:} the
study is \textbf{nocturnal only}, by title, abstract, and design, so it cannot
support a daytime value and is not used for one
(\Cref{sec:derived-archetypes}). \\
\addlinespace
\texttt{cheung2022} & Full &
Open access, IOPscience, read directly. \emph{Caveat recorded:} the
measured-against-modelled gap is specific to agricultural field studies; the
same paper measures \SIrange{-2}{-1}{\degreeCelsius} in irrigated urban parks.
The methodology therefore does not claim that models overstate irrigation
cooling in general. \\
\addlinespace
\texttt{kumar2024} & Full &
PMC open access record, PMC10909648. Two roles: a \emph{validation} of the
inherited calibration (street trees up to \qty{2.8}{\degreeCelsius} in situ,
inside the adopted envelope), and a \emph{citable evidence gap} for productive
landscapes. Its higher figures for other types are deliberately not adopted
(\Cref{sec:non-adoption}). \\
\addlinespace
\texttt{bowler2010} & Metadata + finding (secondary) &
Metadata from the Bangor University repository; the \qty{0.94}{\degreeCelsius}
finding from the publisher abstract description. Elsevier full text blocked. \\
\addlinespace
\texttt{gunawardena2017} & Metadata + finding (secondary) &
Metadata from TU Delft and Bath research portals; findings from the publisher
abstract description. Elsevier full text blocked. Justifies the humidity caveat
applied to blue typologies in tropical-wet climates. \\
\addlinespace
\texttt{venter2021} & Metadata + finding (secondary) &
Metadata from the ADS record and an author-hosted PDF; the surface-versus-canopy
comparison from an indexed description. Publisher full text blocked. Basis for
treating LST as a relative screening proxy only. \\
\addlinespace
\texttt{akbari2001} & Metadata + finding (secondary) &
Multiple indexing records and an LBNL institutional listing. Basis for
\emph{deriving} cooling-energy savings from estimated air temperature
reduction, replacing unsourced per-typology energy factors. \\
\addlinespace
\texttt{yao2023a} & Metadata + finding (secondary) &
Arizona State University Elsevier Pure institutional record, full abstract;
Elsevier full text blocked. The most direct field evidence for large water
bodies, and the reason that archetype's floor was lowered from the
\qty{0.5}{\degreeCelsius} of the approved design to \qty{0.1}{\degreeCelsius}
at implementation (\Cref{sec:water-body}). \\
\addlinespace
\texttt{yao2023b} & Metadata + finding (secondary) &
Arizona State University Elsevier Pure record, full abstract; Elsevier full
text blocked. Independent field corroboration of \texttt{jacobs2020}'s
magnitudes for small water features, and one half of the nocturnal-warming
caveat. \\
\addlinespace
\texttt{gill2013} & Metadata + finding (secondary) &
Crossref record for bibliographic details; abstract via indexing services;
Elsevier full text blocked. \emph{Caveat recorded:} this is a \emph{modelling}
study, cited for the direction and mechanism of drought sensitivity and never
for a numeric cooling value. \\
\addlinespace
\texttt{ouyang2024} & Metadata + finding (secondary) &
Crossref and NASA ADS for bibliographic details; the quantitative findings from
two independent renderings of the publisher abstract; ScienceDirect blocked.
Separates the comfort axis from the air-temperature axis for shade
structures. \\
\addlinespace
\texttt{norton2015} & Metadata + framework description (secondary) &
Metadata from the Monash University research portal; framework description
secondary. The closest published antecedent to this tool. No numeric value is
taken from it. \\
\addlinespace
\texttt{rahman2020} & Metadata only &
Publisher listing. Used qualitatively --- to justify that canopy condition
modulates performance --- never for a numeric value. \\
\addlinespace
\texttt{oke1976} & Metadata only &
Widely indexed. Cited for the conceptual distinction between canopy-layer,
boundary-layer, and surface urban heat islands, not for a numeric value. \\
\end{longtable}

\subsection*{Which tool values depend on non-full verification}

Four dependencies are worth naming explicitly.

The \textbf{energy derivation} of \Cref{sec:energy} rests entirely on
\texttt{akbari2001}, verified at \emph{metadata + finding (secondary)}. The
\SIrange{2}{4}{\percent} per \si{\degreeCelsius} sensitivity is the sole input
to every energy and greenhouse-gas figure the tool produces.

The \textbf{small green area and established park envelopes} depend on
\texttt{bowler2010}'s \qty{0.94}{\degreeCelsius} mean, verified at the same
level, which is also the empirical anchor for the modelling-bias correction of
\Cref{sec:modelling-bias}.

The \textbf{large water body floor} of \qty{0.1}{\degreeCelsius} rests on
\texttt{yao2023a}, also at \emph{metadata + finding (secondary)} --- the
publisher blocks full text, and the
\SIrange{0.1}{0.6}{\degreeCelsius} figure was confirmed from the institutional
record's full abstract. That level of verification was nonetheless sufficient
to establish that the approved design had attributed to this paper a figure it
does not contain, and to move the value (\Cref{sec:water-body}). The episode is
the clearest illustration of what this register is for.

The \textbf{LST proxy restriction} of \Cref{sec:lst} rests on
\texttt{venter2021}, again at the same level --- though here the tool's use is
qualitative, and the conclusion is independently supported by the conceptual
distinction in \texttt{oke1976}.

In each case the tool's confidence ratings and stated limitations reflect the
verification level, and in no case does a value verified only at
\emph{metadata only} carry a number.
